% Chapter 5: Classical Machine Learning Algorithms

\chapter{Classical Machine Learning Algorithms}
\label{chap:classical-ml}

This chapter reviews traditional machine learning methods that provide context and motivation for deep learning approaches. Understanding these classical algorithms helps appreciate the advantages and innovations of deep learning.

\section*{Learning Objectives}
\addcontentsline{toc}{section}{Learning Objectives}

After studying this chapter, you will be able to:

\begin{enumerate}
    \item \textbf{Understand the mathematical foundations} of classical machine learning algorithms including linear regression, logistic regression, and support vector machines
    \item \textbf{Compare and contrast} different approaches to classification and regression problems
    \item \textbf{Implement and optimize} classical algorithms using both closed-form solutions and iterative methods
    \item \textbf{Apply ensemble methods} like random forests and gradient boosting to improve model performance
    \item \textbf{Evaluate the trade-offs} between classical methods and deep learning approaches
    \item \textbf{Choose appropriate algorithms} based on dataset characteristics, computational constraints, and interpretability requirements
    \item \textbf{Understand the limitations} of classical methods that motivated the development of deep learning
    \item \textbf{Apply regularization techniques} to prevent overfitting in classical machine learning models
\end{enumerate}

\difficulty{intermediate} This chapter assumes familiarity with linear algebra, probability theory, and basic optimization concepts from previous chapters.

\input{chapters/chap05-sec01}
\input{chapters/chap05-sec02}
\input{chapters/chap05-sec03}
\input{chapters/chap05-sec04}
\input{chapters/chap05-sec05}
\input{chapters/chap05-sec06}

\section*{Problems}
\addcontentsline{toc}{section}{Problems}

\subsection*{Easy Problems}

\begin{problem}[Linear Regression Basics]
\label{prob:linear-regression-basics}
Given the dataset $\{(1, 2), (2, 4), (3, 6), (4, 8)\}$, find the linear regression model $\hat{y} = wx + b$ that minimizes the mean squared error.

\textbf{Hint:} Use the normal equation $\vect{w}^* = (\mat{X}^\top \mat{X})^{-1} \mat{X}^\top \vect{y}$ where $\mat{X}$ includes a column of ones for the bias term.
\end{problem}

\begin{problem}[Logistic Regression Decision Boundary]
\label{prob:logistic-decision-boundary}
For a logistic regression model with weights $\vect{w} = [2, -1]$ and bias $b = 0$, find the decision boundary equation and classify the point $(1, 1)$.

\textbf{Hint:} The decision boundary occurs where $P(y=1|\vect{x}) = 0.5$, which means $\vect{w}^\top \vect{x} + b = 0$.
\end{problem}

\begin{problem}[SVM Support Vectors]
\label{prob:svm-support-vectors}
Consider a linearly separable dataset with two classes. If you have 3 support vectors, what is the maximum number of training examples you can have?

\textbf{Hint:} Support vectors are the points closest to the decision boundary. Think about the minimum number of support vectors needed for a linearly separable problem.
\end{problem}

\begin{problem}[Decision Tree Splitting]
\label{prob:decision-tree-splitting}
Given a node with 10 examples: 6 belong to class A and 4 to class B, calculate the Gini impurity and entropy.

\textbf{Hint:} Gini impurity = $1 - \sum_{k} p_k^2$ and entropy = $-\sum_{k} p_k \log p_k$ where $p_k$ is the proportion of class $k$.
\end{problem}

\begin{problem}[k-NN Distance Calculation]
\label{prob:knn-distance}
Calculate the Euclidean and Manhattan distances between points $(1, 2, 3)$ and $(4, 5, 6)$.

\textbf{Hint:} Euclidean distance = $\sqrt{\sum_i (x_i - x_i')^2}$ and Manhattan distance = $\sum_i |x_i - x_i'|$.
\end{problem}

\begin{problem}[Regularization Effect]
\label{prob:regularization-effect}
Explain why L1 regularization (Lasso) can drive some weights to exactly zero, while L2 regularization (Ridge) cannot.

\textbf{Hint:} Consider the shape of the L1 and L2 penalty functions and their derivatives at zero.
\end{problem}

\subsection*{Medium Problems}

\begin{problem}[Ridge Regression Derivation]
\label{prob:ridge-derivation}
Derive the closed-form solution for ridge regression: $\vect{w}^* = (\mat{X}^\top \mat{X} + \lambda \mat{I})^{-1} \mat{X}^\top \vect{y}$.

\textbf{Hint:} Start with the ridge regression objective function $L(\vect{w}) = \|\mat{X}\vect{w} - \vect{y}\|^2 + \lambda \|\vect{w}\|^2$ and take the gradient with respect to $\vect{w}$.
\end{problem}

\begin{problem}[SVM Dual Formulation]
\label{prob:svm-dual}
Show that the SVM optimization problem can be written in dual form as:
$$\max_{\boldsymbol{\alpha}} \sum_{i=1}^{n} \alpha_i - \frac{1}{2} \sum_{i=1}^{n} \sum_{j=1}^{n} \alpha_i \alpha_j y^{(i)} y^{(j)} \vect{x}^{(i)} \cdot \vect{x}^{(j)}$$

\textbf{Hint:} Use Lagrange multipliers and the Karush-Kuhn-Tucker (KKT) conditions.
\end{problem}

\begin{problem}[Random Forest Bias-Variance]
\label{prob:random-forest-bias-variance}
Explain how random forests reduce variance compared to a single decision tree. What happens to bias?

\textbf{Hint:} Consider how averaging multiple models affects the bias and variance of the ensemble.
\end{problem}

\begin{problem}[Gradient Boosting Algorithm]
\label{prob:gradient-boosting}
Implement one iteration of gradient boosting for regression. Given current predictions $\hat{y}^{(0)} = [2, 4, 6, 8]$ and true values $y = [3, 5, 7, 9]$, compute the residuals and explain how you would fit the next tree.

\textbf{Hint:} Residuals are $r_i = y_i - \hat{y}_i$. Fit a decision tree to predict these residuals.
\end{problem}

\begin{problem}[k-NN Curse of Dimensionality]
\label{prob:knn-curse-dimensionality}
Explain why k-NN performance degrades in high-dimensional spaces. Provide a mathematical argument.

\textbf{Hint:} Consider how distances between points change as dimensionality increases and the concentration of measure phenomenon.
\end{problem}

\subsection*{Hard Problems}

\begin{problem}[SVM Kernel Trick]
\label{prob:svm-kernel-trick}
Prove that the polynomial kernel $k(\vect{x}, \vect{x}') = (\vect{x}^\top \vect{x}' + c)^d$ is a valid kernel function by showing it corresponds to an inner product in some feature space.

\textbf{Hint:} Expand the polynomial and show it can be written as $\phi(\vect{x})^\top \phi(\vect{x}')$ for some feature map $\phi$.
\end{problem}

\begin{problem}[Logistic Regression Maximum Likelihood]
\label{prob:logistic-mle}
Derive the maximum likelihood estimator for logistic regression. Show that the cross-entropy loss is equivalent to the negative log-likelihood.

\textbf{Hint:} Start with the likelihood function for binary classification and take the negative logarithm.
\end{problem}

\begin{problem}[Decision Tree Pruning]
\label{prob:decision-tree-pruning}
Design an algorithm for decision tree pruning that balances model complexity and performance. What criteria would you use?

\textbf{Hint:} Consider cost-complexity pruning and how to measure the trade-off between tree size and accuracy.
\end{problem}

\begin{problem}[Ensemble Methods Theory]
\label{prob:ensemble-theory}
Prove that for an ensemble of $B$ independent models with error rate $p < 0.5$, the ensemble error rate approaches 0 as $B \to \infty$.

\textbf{Hint:} Use the binomial distribution and the fact that the ensemble makes an error only when more than half of the models are wrong.
\end{problem}

\begin{problem}[Feature Selection for k-NN]
\label{prob:feature-selection-knn}
Design a feature selection algorithm specifically for k-NN that addresses the curse of dimensionality. How would you evaluate feature importance?

\textbf{Hint:} Consider wrapper methods, filter methods, and embedded methods. Think about how to measure the contribution of each feature to the k-NN decision.
\end{problem}
